\documentclass{amsart}
\usepackage{amsmath,amssymb,amsthm,hyperref}

\newtheorem{theorem}{Theorem}
\newtheorem{lemma}{Lemma}
\newtheorem{proposition}{Proposition}
\newtheorem{corollary}{Corollary}
\theoremstyle{definition}
\newtheorem{definition}{Definition}
\newtheorem{remark}{Remark}

\newcommand{\Ord}{\mathrm{Ord}}
\newcommand{\interp}[1]{[\![#1]\!]}
\newcommand{\To}{\to_\beta}
\newcommand{\FV}{\mathrm{FV}}

\begin{document}

\title{An elementary ordinal assignment witnessing strong normalization of STLC}
\maketitle

\section{Statement}

Let $\Lambda_\to$ be the set of simply typed $\lambda$-terms over a set of base
types, with $\beta$-reduction $\To$. We construct a map
\[
  \# : \Lambda_\to \longrightarrow \Ord
\]
such that
\[
  M \To N \;\Longrightarrow\; \#M > \#N \qquad (M,N\in\Lambda_\to),
\]
where the construction of $\#$ is \emph{elementary}: it is defined by structural
recursion on types and terms and never invokes strong normalization (SN). Since
$\Ord$ is well-founded, the existence of $\#$ yields at once a direct proof of SN:
an infinite reduction $M_0\To M_1\To\cdots$ would give an infinite strictly
descending sequence $\#M_0 > \#M_1 > \cdots$ of ordinals, which is impossible.

The map $\#$ is in fact \emph{natural-number} valued (we land in $\omega\subseteq\Ord$),
which is more than enough; for pure STLC the longest reduction sequence is finite, so
ordinals below $\omega$ suffice. We phrase the codomain as $\Ord$ to match the problem
statement and to make the well-foundedness step explicit.

\section{Conventions on the calculus}

\emph{Simple types} are generated by $\sigma ::= o \mid (\sigma\to\tau)$, where $o$
ranges over a fixed set of base types (the argument treats all base types identically,
so we may assume a single base type $o$). \emph{Terms} are typed as usual; we write
$M:\sigma$ when $M$ has type $\sigma$, work up to $\alpha$-equivalence, write $\FV(M)$
for the free variables, and write $M[x:=N]$ for capture-avoiding substitution. The
reduction relation $\To$ is the \emph{full} compatible (congruence) closure of
\[
  (\lambda x.M)\,N \;\To\; M[x:=N];
\]
a redex may be contracted anywhere inside a term. Nothing below presupposes that $\To$
is normalizing.

\section{The interpretation domains}

For each type $\sigma$ we define a set $\mathcal M_\sigma$ of \emph{functionals}, two
relations $<_\sigma$ (strict) and $\le_\sigma$ (non-strict) on $\mathcal M_\sigma$, and
a \emph{weight} $\|\cdot\|_\sigma:\mathcal M_\sigma\to\mathbb N_{\ge 1}$, all by
recursion on $\sigma$. Write $\mathbb P=\{1,2,3,\dots\}$ with the usual order.

\begin{definition}[Domains, orders, weight]\label{def:dom}
\leavevmode
\begin{itemize}
\item \textbf{Base type.}
  $\mathcal M_o=\mathbb P$, with $m<_o n :\Leftrightarrow m<n$,
  $m\le_o n:\Leftrightarrow m\le n$, and $\|n\|_o:=n$.
\item \textbf{Arrow type.}
  $\mathcal M_{\sigma\to\tau}$ is the set of all functions
  $f:\mathcal M_\sigma\to\mathcal M_\tau$ that are \emph{simultaneously strictly and
  non-strictly monotone}:
  \begin{align*}
    \text{(S)}&\quad a<_\sigma a' \;\Longrightarrow\; f(a)<_\tau f(a')
       &&(a,a'\in\mathcal M_\sigma),\\
    \text{(N)}&\quad a\le_\sigma a' \;\Longrightarrow\; f(a)\le_\tau f(a')
       &&(a,a'\in\mathcal M_\sigma).
  \end{align*}
  They are ordered everywhere-strict and everywhere-non-strict \emph{pointwise}:
  \[
    f<_{\sigma\to\tau}g :\Leftrightarrow
    \forall a\in\mathcal M_\sigma,\ f(a)<_\tau g(a),
    \qquad
    f\le_{\sigma\to\tau}g :\Leftrightarrow
    \forall a\in\mathcal M_\sigma,\ f(a)\le_\tau g(a),
  \]
  with weight read off at a canonical least argument:
  \[
    \|f\|_{\sigma\to\tau}:=\|f(\mathbf 0_\sigma)\|_\tau .
  \]
\end{itemize}
The \emph{canonical element} $\mathbf 0_\sigma$ and the \emph{lifts}
$\mathrm{lift}_\sigma:\mathbb P\to\mathcal M_\sigma$ are defined by
\[
  \mathrm{lift}_o(k):=k,\qquad
  \mathrm{lift}_{\sigma\to\tau}(k):=\bigl(a\mapsto\mathrm{lift}_\tau(k+\|a\|_\sigma)\bigr),
  \qquad \mathbf 0_\sigma:=\mathrm{lift}_\sigma(1).
\]
\end{definition}

\begin{remark}[Why both (S) and (N), and why $\le_\sigma$ is pointwise]\label{rem:leq}
Two design choices are essential and interlocking.

\emph{(i) $\le_\sigma$ is defined pointwise (recursively), not as ``$<_\sigma$ or $=$''.}
At an arrow type these genuinely differ: if $f,g$ are monotone with $f(a)\le_\tau g(a)$ for
all $a$, yet $f(a_0)=g(a_0)$ for some $a_0$ while $f(a_1)<_\tau g(a_1)$ for some other $a_1$,
then $f\le_{\sigma\to\tau}g$ holds while neither $f<_{\sigma\to\tau}g$ nor $f=g$ holds. Such
configurations arise in the monotonicity argument; the pointwise $\le_\sigma$ is the right
inductive invariant.

\emph{(ii) We require both (S) \emph{and} (N) for membership.} A strictly monotone functional
need \emph{not} be non-strictly monotone in the pointwise sense of~(i). Concretely, take the
base codomain $\mathcal M_o=\mathbb P$ and two functions $a=(2,4)$, $a'=(3,4)$ in
$\mathcal M_{o\to o}$ (writing a domain-$\{1,2\}$ table); then $a\le_{o\to o}a'$ but not
$a<_{o\to o}a'$. One can build a functional $F:\mathcal M_{o\to o}\to\mathcal M_o$ that is
strictly monotone for $<$ (i.e.\ $u<v$ everywhere $\Rightarrow F(u)<F(v)$) yet has
$F(a)>F(a')$, so $F$ violates~(N). \emph{(We verified the existence of such an $F$ by a finite
search; see the verification note.)} Hence non-strict monotonicity is an independent property
that must be \emph{maintained as part of membership}; merely lying in a space of strictly
monotone functions does not give it for free. Building (N) into Definition~\ref{def:dom} is
exactly what makes the application case of the interpretation (where we apply one functional
to another) preserve non-strict monotonicity automatically, with no auxiliary construction.
\end{remark}

\begin{lemma}[Basic properties]\label{lem:basic}
For every type $\sigma$:
\begin{enumerate}
\item[(a)] For each $k\in\mathbb P$, $\mathrm{lift}_\sigma(k)\in\mathcal M_\sigma$, and
  $\mathrm{lift}_\sigma$ is monotone in its integer argument in both senses:
  $j<j'\Rightarrow\mathrm{lift}_\sigma(j)<_\sigma\mathrm{lift}_\sigma(j')$ and
  $j\le j'\Rightarrow\mathrm{lift}_\sigma(j)\le_\sigma\mathrm{lift}_\sigma(j')$.
  In particular $\mathbf 0_\sigma\in\mathcal M_\sigma$.
\item[(b)] The weight is order-preserving in both senses:
  $v<_\sigma v'\Rightarrow\|v\|_\sigma<\|v'\|_\sigma$ and
  $v\le_\sigma v'\Rightarrow\|v\|_\sigma\le\|v'\|_\sigma$.
\item[(c)] $\le_\sigma$ is a partial order; $<_\sigma$ is a strict partial order; and they
  interleave: $v<_\sigma v'\Rightarrow v\le_\sigma v'$; if $v\le_\sigma v'$ and
  $v'<_\sigma v''$ then $v<_\sigma v''$; if $v<_\sigma v'$ and $v'\le_\sigma v''$ then
  $v<_\sigma v''$.
\item[(d)] $<_\sigma$ is well-founded.
\end{enumerate}
\end{lemma}

\begin{proof}
We prove (a), (b), (c) by simultaneous induction on $\sigma$; (d) follows from (b).

\emph{(c).} At base type these are the standard facts about $(\mathbb P,<,\le)$. At arrow
type every assertion is the universally quantified pointwise statement, which holds by the
IH for $\tau$ applied at each argument $a$. (For antisymmetry: $f\le g$ and $g\le f$ mean
$f(a)\le g(a)\le f(a)$, so $f(a)=g(a)$ for all $a$, i.e.\ $f=g$.)

\emph{(b).} Base: $\|n\|_o=n$. Arrow: if $f<_{\sigma\to\tau}g$ then
$f(\mathbf 0_\sigma)<_\tau g(\mathbf 0_\sigma)$, so by IH(b),
$\|f\|=\|f(\mathbf 0_\sigma)\|_\tau<\|g(\mathbf 0_\sigma)\|_\tau=\|g\|$; the non-strict case
is identical with $\le$. (Here $\mathbf 0_\sigma\in\mathcal M_\sigma$ by IH(a).)

\emph{(a).} Base: $\mathrm{lift}_o(k)=k$, and $j<j'\Rightarrow j<j'$, $j\le j'\Rightarrow
j\le j'$. Arrow $\sigma\to\tau$: first, $\mathrm{lift}_{\sigma\to\tau}(k)=
(a\mapsto\mathrm{lift}_\tau(k+\|a\|_\sigma))$ lies in $\mathcal M_{\sigma\to\tau}$, i.e.\
satisfies (S) and (N). For (S): $a<_\sigma a'\Rightarrow\|a\|_\sigma<\|a'\|_\sigma$ (IH(b))
$\Rightarrow k+\|a\|_\sigma<k+\|a'\|_\sigma\Rightarrow
\mathrm{lift}_\tau(k+\|a\|_\sigma)<_\tau\mathrm{lift}_\tau(k+\|a'\|_\sigma)$ by IH(a) for
$\tau$ (integer-strict-monotonicity). For (N), replace $<$ by $\le$ throughout, using IH(b)
non-strict and IH(a) integer-non-strict-monotonicity for $\tau$. Hence
$\mathrm{lift}_{\sigma\to\tau}(k)\in\mathcal M_{\sigma\to\tau}$. Its integer-monotonicity:
for $j<j'$ and every $a$, $j+\|a\|_\sigma<j'+\|a\|_\sigma$, so
$\mathrm{lift}_\tau(j+\|a\|_\sigma)<_\tau\mathrm{lift}_\tau(j'+\|a\|_\sigma)$ by IH(a) for
$\tau$; this holds for all $a$, hence $\mathrm{lift}_{\sigma\to\tau}(j)<_{\sigma\to\tau}
\mathrm{lift}_{\sigma\to\tau}(j')$. The non-strict integer-monotonicity is the same with
$\le$.

\emph{(d).} By (b), $\|\cdot\|_\sigma$ maps $(\mathcal M_\sigma,<_\sigma)$
strictly-order-preservingly into the well-founded $(\mathbb P,<)$, so an infinite
$<_\sigma$-descending chain would yield an infinite $<$-descending chain in $\mathbb P$;
hence $<_\sigma$ is well-founded.
\end{proof}

\begin{remark}
Lemma~\ref{lem:basic}(d) is the crucial point flagged in the literature: the general
pointwise order on a function space need not be well-founded. We sidestep this by carrying
the explicit weight $\|\cdot\|_\sigma$, which embeds each $(\mathcal M_\sigma,<_\sigma)$
strictly-order-preservingly into the well-founded $(\mathbb P,<)$. All comparisons we
ultimately make descend to comparisons of natural numbers (Theorem~\ref{thm:main}), with no
reference to reduction or SN.
\end{remark}

\begin{definition}[Weight shift]\label{def:add}
For $k\in\mathbb N$ define $\mathrm{add}_\sigma(\cdot,k):\mathcal M_\sigma\to\mathcal M_\sigma$ by
\[
  \mathrm{add}_o(n,k):=n+k,\qquad
  \mathrm{add}_{\sigma\to\tau}(f,k):=\bigl(a\mapsto\mathrm{add}_\tau(f(a),k)\bigr).
\]
\end{definition}

\begin{lemma}[Properties of $\mathrm{add}$]\label{lem:add}
For every type $\sigma$, all $v,v'\in\mathcal M_\sigma$, $k,k'\in\mathbb N$:
\begin{enumerate}
\item[(a)] $\mathrm{add}_\sigma(v,k)\in\mathcal M_\sigma$;
\item[(b)] $\|\mathrm{add}_\sigma(v,k)\|_\sigma=\|v\|_\sigma+k$;
\item[(c)] $v<_\sigma v'\Rightarrow\mathrm{add}_\sigma(v,k)<_\sigma\mathrm{add}_\sigma(v',k)$,
  and $v\le_\sigma v'\Rightarrow\mathrm{add}_\sigma(v,k)\le_\sigma\mathrm{add}_\sigma(v',k)$;
\item[(d)] $k<k'\Rightarrow\mathrm{add}_\sigma(v,k)<_\sigma\mathrm{add}_\sigma(v,k')$,
  and $k\le k'\Rightarrow\mathrm{add}_\sigma(v,k)\le_\sigma\mathrm{add}_\sigma(v,k')$;
  and $\mathrm{add}_\sigma(v,0)=v$.
\end{enumerate}
\end{lemma}

\begin{proof}
Induction on $\sigma$. Base: (a) trivial; (b) $\|n+k\|=n+k$; (c) $n<n'\Rightarrow n+k<n'+k$
and $n\le n'\Rightarrow n+k\le n'+k$; (d) $k<k'\Rightarrow n+k<n+k'$, $k\le k'\Rightarrow
n+k\le n+k'$, $n+0=n$. The non-strict shift clause $k\le k'\Rightarrow\mathrm{add}(v,k)\le
\mathrm{add}(v,k')$ at base type is immediate, and at arrow types it is the pointwise
statement obtained by applying the base case at each argument (or, equivalently, from the
strict clause via $<\Rightarrow\le$ when $k<k'$ and reflexivity when $k=k'$).

Arrow $\sigma\to\tau$. (a) We must check $\mathrm{add}(f,k)$ satisfies (S) and (N). For (S):
$a<_\sigma a'\Rightarrow f(a)<_\tau f(a')\Rightarrow_{\text{IH(c)}}
\mathrm{add}_\tau(f(a),k)<_\tau\mathrm{add}_\tau(f(a'),k)$. For (N): $a\le_\sigma a'\Rightarrow
f(a)\le_\tau f(a')\Rightarrow_{\text{IH(c)}}\mathrm{add}_\tau(f(a),k)\le_\tau
\mathrm{add}_\tau(f(a'),k)$. (Here $f\in\mathcal M_{\sigma\to\tau}$ supplies both
$f(a)<_\tau f(a')$ and $f(a)\le_\tau f(a')$ from its membership clauses (S),(N).) Hence
$\mathrm{add}(f,k)\in\mathcal M_{\sigma\to\tau}$.
(b) $\|\mathrm{add}(f,k)\|=\|\mathrm{add}_\tau(f(\mathbf 0_\sigma),k)\|_\tau=
\|f(\mathbf 0_\sigma)\|_\tau+k=\|f\|+k$ by IH(b). (c),(d) are the pointwise statements via
IH(c),(d); e.g.\ for (c) non-strict, $f\le g$ means $f(a)\le g(a)$ for all $a$, so by IH(c)
$\mathrm{add}_\tau(f(a),k)\le\mathrm{add}_\tau(g(a),k)$ for all $a$, i.e.\
$\mathrm{add}(f,k)\le\mathrm{add}(g,k)$.
\end{proof}

\section{Interpretation of terms}

An \emph{environment} $\rho$ assigns to each variable $x:\tau$ a value
$\rho(x)\in\mathcal M_\tau$; $\rho[x:=a]$ denotes update.

\begin{definition}[Term interpretation]\label{def:interp}
For $M:\sigma$ and an environment $\rho$ on $\FV(M)$, define $\interp{M}_\rho\in\mathcal M_\sigma$
by recursion on $M$:
\begin{align*}
  \interp{x}_\rho &:= \rho(x),\\
  \interp{M\,N}_\rho &:= \interp{M}_\rho\bigl(\interp{N}_\rho\bigr),\\
  \interp{\lambda x^\sigma.M}_\rho &:=
     \Bigl(a\in\mathcal M_\sigma\ \mapsto\
       \mathrm{add}_\tau\!\bigl(\interp{M}_{\rho[x:=a]},\ 1+\|a\|_\sigma\bigr)\Bigr),
     \quad M:\tau .
\end{align*}
\end{definition}

This is plain structural recursion, manifestly total, using no normalization hypothesis.
The decisive design choice is at abstraction: the result is shifted by $1+\|a\|_\sigma$, an
amount that \emph{strictly increases with the argument $a$}. This makes abstractions
strictly monotone \emph{even when $x$ does not occur in $M$}, and furnishes the strict drop
at a redex.

\begin{remark}[Why $1+\|a\|$ and not a fixed ``$+1$'']\label{rem:why}
A naive choice $\interp{\lambda x.M}_\rho=(a\mapsto\interp{M}_{\rho[x:=a]})$, or with a
constant ``$+1$'', \emph{fails}: if $x\notin\FV(M)$ then $\interp{M}_{\rho[x:=a]}$ is
independent of $a$, so the function is constant, hence not strictly monotone, so it does not
lie in $\mathcal M_{\sigma\to\tau}$. This breaks Lemma~\ref{lem:cong}(ii). Shifting by
$1+\|a\|_\sigma$ makes the dependence on $a$ strict, repairing the gap while still leaving a
strict ``$+1$'' for the redex drop.
\end{remark}

\section{The monotonicity lemma}

This is the technical heart. We prove well-definedness and both environment-monotonicities by
one structural induction. Crucially, because non-strict monotonicity in the argument is now
\emph{built into membership} in $\mathcal M_\sigma$ (clause (N) of Definition~\ref{def:dom}),
we never need a separate ``argument-monotonicity'' clause in this induction nor any auxiliary
term: it suffices to show every interpretation \emph{lies in} $\mathcal M_\sigma$, and (N)
then comes along for free.

We restrict environments to assign values in $\mathcal M_\tau$ (which, at arrow types, are by
definition both strictly and non-strictly monotone). This class is closed under everything we
do: the canonical environment uses lifts ($\in\mathcal M$ by Lemma~\ref{lem:basic}(a)), and
every value we ever insert into an environment is either an interpretation $\interp{N}_\theta$
(which lies in $\mathcal M$ by clause (M0) below) or an abstraction's bound argument $a$,
already drawn from $\mathcal M_\sigma$.

\begin{lemma}[Monotonicity of interpretations]\label{lem:official}
For every term $M:\sigma$ and all environments $\rho,\rho'$ on $\FV(M)$ \emph{with values in
the respective $\mathcal M_\tau$}:
\begin{enumerate}
\item[\textup{(M0)}] $\interp{M}_\rho\in\mathcal M_\sigma$;
\item[\textup{(M$\le$)}] if $\rho(y)\le_{\tau_y}\rho'(y)$ for all $y\in\FV(M)$, then
  $\interp{M}_\rho\le_\sigma\interp{M}_{\rho'}$;
\item[\textup{(M$<$)}] if moreover $\rho(x)<_{\tau_x}\rho'(x)$ for some $x\in\FV(M)$ and
  $\rho(y)=\rho'(y)$ for $y\ne x$, then $\interp{M}_\rho<_\sigma\interp{M}_{\rho'}$.
\end{enumerate}
\textup{(}As an immediate consequence of \textup{(M0)} and clause \textup{(N)} of
Definition~\ref{def:dom}: if $\sigma=\tau'\to\sigma''$ then $\interp{M}_\rho$ is non-strictly
monotone in its argument.\textup{)}
\end{lemma}

\begin{proof}
By induction on $M$. We use Lemma~\ref{lem:basic}(c) (order facts, including
$<\Rightarrow\le$ and the mixed-transitivity rules) and Lemma~\ref{lem:add}(c) (that
$\mathrm{add}(\cdot,k)$ preserves $\le$ and $<$). Throughout we freely use that any value of
arrow type appearing in our environments, or produced as an interpretation, lies in
$\mathcal M$ and hence satisfies both (S) and (N).

\medskip\noindent\emph{Variable $M=y$.}
(M0): $\rho(y)\in\mathcal M_{\tau_y}$. (M$\le$): $\rho(y)\le\rho'(y)$. (M$<$): $x=y$ gives
$\rho(y)<\rho'(y)$.

\medskip\noindent\emph{Abstraction $M=\lambda z^{\tau'}.P$, $P:\sigma'$,
$\sigma=\tau'\to\sigma'$.} Rename $z$ fresh. For any argument $a\in\mathcal M_{\tau'}$, the
updated environment $\rho[z:=a]$ still has all values in the relevant $\mathcal M$ ($a\in
\mathcal M_{\tau'}$ and the other values inherited from $\rho$), so the IH applies to $P$ at
such environments.
\begin{itemize}
\item (M0). We show $\interp{M}_\rho:=\bigl(a\mapsto\mathrm{add}(\interp{P}_{\rho[z:=a]},
  1{+}\|a\|)\bigr)$ satisfies both (S) and (N), hence lies in $\mathcal M_{\tau'\to\sigma'}$.

  \emph{(S).} Fix $a<_{\tau'}a'$. By IH \textup{(M$\le$)} for $P$ (environments differ only at
  $z$, where $a\le a'$ by $<\Rightarrow\le$),
  $\interp{P}_{\rho[z:=a]}\le_{\sigma'}\interp{P}_{\rho[z:=a']}$. By Lemma~\ref{lem:basic}(b),
  $\|a\|_{\tau'}<\|a'\|_{\tau'}$, hence $1{+}\|a\|<1{+}\|a'\|$. Then by
  Lemma~\ref{lem:add}(c),(d) and mixed transitivity,
  \[
    \interp{M}_\rho(a)=\mathrm{add}(\interp{P}_{\rho[z:=a]},1{+}\|a\|)
    \le\mathrm{add}(\interp{P}_{\rho[z:=a']},1{+}\|a\|)
    <\mathrm{add}(\interp{P}_{\rho[z:=a']},1{+}\|a'\|)=\interp{M}_\rho(a').
  \]

  \emph{(N).} Fix $a\le_{\tau'}a'$. By IH \textup{(M$\le$)} for $P$,
  $\interp{P}_{\rho[z:=a]}\le_{\sigma'}\interp{P}_{\rho[z:=a']}$; and $\|a\|\le\|a'\|$
  (Lemma~\ref{lem:basic}(b)), so $1{+}\|a\|\le 1{+}\|a'\|$. Then by Lemma~\ref{lem:add}(c)
  (value side, $\le$) and the shift-side non-strict monotonicity
  $k\le k'\Rightarrow\mathrm{add}(v,k)\le\mathrm{add}(v,k')$ (Lemma~\ref{lem:add}(d)),
  composed by transitivity of $\le$,
  \[
    \interp{M}_\rho(a)=\mathrm{add}(\interp{P}_{\rho[z:=a]},1{+}\|a\|)
    \le_{\sigma'}\mathrm{add}(\interp{P}_{\rho[z:=a']},1{+}\|a'\|)=\interp{M}_\rho(a').
  \]
  Thus $\interp{M}_\rho\in\mathcal M_{\tau'\to\sigma'}$.
\item (M$\le$)/(M$<$). For each $a$, apply IH \textup{(M$\le$)}/\textup{(M$<$)} for $P$ at
  $\rho[z:=a],\rho'[z:=a]$, then $\mathrm{add}(\cdot,1{+}\|a\|)$ (Lemma~\ref{lem:add}(c)).
  Holding for all $a$, this gives $\interp{M}_\rho\le_{\tau'\to\sigma'}\interp{M}_{\rho'}$,
  resp.\ $<$, by the pointwise definitions (Definition~\ref{def:dom}).
\end{itemize}

\medskip\noindent\emph{Application $M=PQ$, $P:\tau'\to\sigma$, $Q:\tau'$.}
Let $p=\interp{P}_\rho,p'=\interp{P}_{\rho'},q=\interp{Q}_\rho,q'=\interp{Q}_{\rho'}$.

(M0): by IH \textup{(M0)} for $P,Q$, $p,p'\in\mathcal M_{\tau'\to\sigma}$ and
$q,q'\in\mathcal M_{\tau'}$. Since $p\in\mathcal M_{\tau'\to\sigma}$ is a function
$\mathcal M_{\tau'}\to\mathcal M_\sigma$, we have $\interp{PQ}_\rho=p(q)\in\mathcal M_\sigma$.
In particular, if $\sigma$ is an arrow type, $p(q)\in\mathcal M_\sigma$ \emph{already}
satisfies (N) (non-strict argument monotonicity), \emph{by definition of $\mathcal M_\sigma$};
no separate argument is needed.

For (M$\le$) and (M$<$) we use the membership $p'\in\mathcal M_{\tau'\to\sigma}$, whose
clause (N) gives that $p'$ is non-strictly monotone in its argument:
\begin{equation}\label{eq:nmP}
  q\le_{\tau'}q'\ \Longrightarrow\ p'(q)\le_\sigma p'(q').
\end{equation}
\begin{itemize}
\item (M$\le$). By IH \textup{(M$\le$)} for $P$, $p\le_{\tau'\to\sigma}p'$, so pointwise at
  $q$, $p(q)\le_\sigma p'(q)$. By IH \textup{(M$\le$)} for $Q$, $q\le_{\tau'}q'$, so by
  \eqref{eq:nmP}, $p'(q)\le_\sigma p'(q')$. Transitivity of $\le$ gives
  $p(q)\le_\sigma p'(q')$, i.e.\ $\interp{PQ}_\rho\le_\sigma\interp{PQ}_{\rho'}$.
\item (M$<$). Either $x\in\FV(Q)$ or $x\in\FV(P)$ (with $\rho=\rho'$ off $x$).
  \begin{itemize}
  \item If $x\in\FV(Q)$: by IH \textup{(M$<$)} for $Q$, $q<_{\tau'}q'$; since $p'$ satisfies
    (S) (it lies in $\mathcal M_{\tau'\to\sigma}$), $p'(q)<_\sigma p'(q')$; and $p\le p'$
    gives $p(q)\le_\sigma p'(q)$; mixed transitivity ($\le$ then $<$) gives
    $p(q)<_\sigma p'(q')$.
  \item If $x\notin\FV(Q)$ and $x\in\FV(P)$: then $q=q'$, and IH \textup{(M$<$)} for $P$
    gives $p<_{\tau'\to\sigma}p'$, so pointwise at $q$, $p(q)<_\sigma p'(q)=p'(q')$.
  \end{itemize}
  Either way $\interp{PQ}_\rho<_\sigma\interp{PQ}_{\rho'}$.
\end{itemize}
This completes the induction.
\end{proof}

\begin{remark}[The key simplification]\label{rem:keyfix}
Non-strict argument-monotonicity of an interpreted functional is \emph{not} the (false)
statement that every strictly monotone functional is non-strictly monotone
(Remark~\ref{rem:leq}(ii)). It is secured by \emph{carrying it as part of membership} in
$\mathcal M$ (clause (N) of Definition~\ref{def:dom}) and \emph{verifying it is preserved}
by every operation: by lifts and $\mathrm{add}$ (Lemmas~\ref{lem:basic}(a),
\ref{lem:add}(a)) and by the two term constructors. For abstraction this is the explicit
(N) bullet of the (M0) case above; for application it is automatic, since $p(q)$ is a member
of the codomain $\mathcal M_\sigma$, which by definition contains only functionals satisfying
(N). This avoids any appeal to a non-subterm auxiliary, keeping the induction strictly
structural on $M$.
\end{remark}

\begin{remark}[The three difficulties, discharged]\label{rem:threehard}
We record explicitly where Lemma~\ref{lem:official} settles the three points that make
strict (rather than merely weak) monotonicity delicate.
\emph{(i) Strictness, not weak monotonicity.} Two \emph{distinct} strict statements are
proved. Strictness \emph{in a free variable} is clause (M$<$); strictness \emph{in the
abstraction's own argument} is clause (S) of (M0), established in the abstraction case by the
increment $1{+}\|a\|_{\tau'}$ together with Lemma~\ref{lem:add}(c),(d) and mixed transitivity
(Lemma~\ref{lem:basic}(c)) --- and it holds \emph{even when} $z\notin\FV(P)$, where the
underlying $\interp{P}_{\rho[z:=a]}$ is constant in $a$. This argument-strictness is exactly
what drives the reduction-in-argument case via Lemma~\ref{lem:cong}(ii).
\emph{(ii) Propagation through application.} In (M$<$), application, when the strictly
increased variable lies in the argument $Q$ we use that $p'=\interp{P}_{\rho'}\in
\mathcal M_{\tau'\to\sigma}$ satisfies (S), so $q<_{\tau'}q'\Rightarrow p'(q)<_\sigma p'(q')$;
when it lies in the function $P$ we use the pointwise $p<_{\tau'\to\sigma}p'$ at $q=q'$. In
both sub-cases the strict drop survives functional application.
\emph{(iii) Compatibility of the increment with the closure conditions.} The increment is the
map $\mathrm{add}_{\sigma'}(\cdot,1{+}\|a\|)$, which preserves membership
(Lemma~\ref{lem:add}(a)) and both orders (Lemma~\ref{lem:add}(c),(d)); the (M0) abstraction
case then verifies \emph{both} (S) and (N) for the resulting functional, so the incremented
abstraction lies in $\mathcal M_{\tau'\to\sigma'}$ and remains strictly monotone. Thus the
increment is fully compatible with the defining closure conditions (S),(N) of
Definition~\ref{def:dom}.
\end{remark}

\section{Substitution, redex drop, congruence, and the main theorem}

\begin{lemma}[Substitution]\label{lem:subst}
Let $M:\sigma$, $x:\tau$, $N:\tau$, and let $\rho$ be an environment on $\FV(M[x:=N])$ with
values in the respective $\mathcal M$. Then
$\interp{M[x:=N]}_\rho=\interp{M}_{\rho[x:=\interp{N}_\rho]}$.
\end{lemma}

\begin{proof}
Induction on $M$, with $\alpha$-conventions so bound variables of $M$ avoid $x$ and
$\FV(N)$. We use that $\interp{T}_\rho$ depends only on $\rho{\restriction}\FV(T)$. Note the
inserted value $\interp{N}_\rho$ lies in $\mathcal M_\tau$ (Lemma~\ref{lem:official},
(M0)), so $\rho[x:=\interp N_\rho]$ respects the environment discipline.
\begin{itemize}
\item $M=x$: both sides $=\interp N_\rho$. \quad $M=y\neq x$: both sides $=\rho(y)$.
\item $M=PQ$: by IH, $\interp{(PQ)[x:=N]}_\rho=\interp{P[x:=N]}_\rho(\interp{Q[x:=N]}_\rho)
  =\interp{P}_{\rho[x:=\interp N_\rho]}(\interp{Q}_{\rho[x:=\interp N_\rho]})
  =\interp{PQ}_{\rho[x:=\interp N_\rho]}$.
\item $M=\lambda z.P$, $z\notin\{x\}\cup\FV(N)$. For all $a$,
  \begin{align*}
   \interp{(\lambda z.P)[x:=N]}_\rho(a)
   &=\mathrm{add}\bigl(\interp{P[x:=N]}_{\rho[z:=a]},1{+}\|a\|\bigr)\\
   &\overset{\text{IH}}{=}\mathrm{add}\bigl(\interp{P}_{\rho[z:=a][x:=\interp N_{\rho[z:=a]}]},1{+}\|a\|\bigr)\\
   &=\mathrm{add}\bigl(\interp{P}_{\rho[x:=\interp N_\rho][z:=a]},1{+}\|a\|\bigr)
   =\interp{\lambda z.P}_{\rho[x:=\interp N_\rho]}(a),
  \end{align*}
  using $\interp N_{\rho[z:=a]}=\interp N_\rho$ ($z\notin\FV(N)$) and that updates at
  $z\neq x$ commute. Equality for all $a$ gives equality of the functions.
\end{itemize}
\end{proof}

\begin{lemma}[Strict drop at a redex]\label{lem:redex}
For every redex $(\lambda x^\tau.M)\,N$ of type $\sigma$ and every environment $\rho$ on its
free variables (with values in the respective $\mathcal M$),
$\interp{M[x:=N]}_\rho<_\sigma\interp{(\lambda x^\tau.M)\,N}_\rho$.
\end{lemma}

\begin{proof}
Let $a:=\interp N_\rho\in\mathcal M_\tau$; then $\|a\|_\tau\ge1$, so $1+\|a\|_\tau>0$. By
Definition~\ref{def:interp},
$\interp{(\lambda x.M)N}_\rho=\interp{\lambda x.M}_\rho(a)=
\mathrm{add}_\sigma(\interp{M}_{\rho[x:=a]},1{+}\|a\|_\tau)$, and by Lemma~\ref{lem:subst},
$\interp{M[x:=N]}_\rho=\interp{M}_{\rho[x:=a]}$. By Lemma~\ref{lem:add}(d),
\[
  \interp{M[x:=N]}_\rho=\mathrm{add}_\sigma(\interp{M}_{\rho[x:=a]},0)
  <_\sigma\mathrm{add}_\sigma(\interp{M}_{\rho[x:=a]},1{+}\|a\|_\tau)=\interp{(\lambda x.M)N}_\rho.\qedhere
\]
\end{proof}

\begin{remark}[Robustness]
The drop comes entirely from the strictly positive shift at the abstraction.
\emph{Discarding} ($x\notin\FV(M)$): $\interp{M}_{\rho[x:=a]}=\interp{M}_\rho$ is independent
of $N$, yet $\mathrm{add}(\interp M_\rho,0)<\mathrm{add}(\interp M_\rho,1{+}\|a\|)$.
\emph{Duplicating} ($x$ occurs $\ge2$ times): the Substitution Lemma uses the single value
$a=\interp N_\rho$ wherever $x$ occurs, so no copy of $N$ is re-evaluated and the measure
cannot blow up.
\end{remark}

\begin{lemma}[Congruence]\label{lem:cong}
Suppose $M,M':\sigma$ satisfy $\interp{M}_\rho<_\sigma\interp{M'}_\rho$ for every
\emph{(}admissible\emph{)} environment $\rho$. Then for every such $\rho$:
\textup{(i)} $\interp{M\,P}_\rho<\interp{M'\,P}_\rho$;
\textup{(ii)} $\interp{P\,M}_\rho<\interp{P\,M'}_\rho$;
\textup{(iii)} $\interp{\lambda y.M}_\rho<\interp{\lambda y.M'}_\rho$.
\end{lemma}

\begin{proof}
\emph{(i)} $M,M':\tau'\to\sigma$, $P:\tau'$. The hypothesis is
$\interp{M}_\rho(b)<_\sigma\interp{M'}_\rho(b)$ for all $b$; take $b=\interp P_\rho$.

\emph{(ii)} $P:\tau'\to\sigma$, $M,M':\tau'$. By hypothesis
$\interp{M}_\rho<_{\tau'}\interp{M'}_\rho$. Since $\interp P_\rho\in\mathcal M_{\tau'\to\sigma}$
satisfies (S) (Lemma~\ref{lem:official}, (M0)),
$\interp{PM}_\rho=\interp P_\rho(\interp M_\rho)<_\sigma\interp P_\rho(\interp{M'}_\rho)=\interp{PM'}_\rho$.

\emph{(iii)} For $y:\tau'$ and every $a$, the hypothesis at $\rho[y:=a]$ gives
$\interp{M}_{\rho[y:=a]}<_\sigma\interp{M'}_{\rho[y:=a]}$, and Lemma~\ref{lem:add}(c) yields
$\interp{\lambda y.M}_\rho(a)<_\sigma\interp{\lambda y.M'}_\rho(a)$ for all $a$, hence
$\interp{\lambda y.M}_\rho<_{\tau'\to\sigma}\interp{\lambda y.M'}_\rho$.
\end{proof}

\begin{proposition}[Strict decrease under one step]\label{prop:step}
If $M\To N$ then $M,N:\sigma$ for a common $\sigma$, and for every admissible environment
$\rho$ on $\FV(M)\supseteq\FV(N)$, $\interp{N}_\rho<_\sigma\interp{M}_\rho$.
\end{proposition}

\begin{proof}
Induction on the context of the contracted redex.
\emph{Head:} $M=(\lambda x.P)Q\To P[x:=Q]=N$ -- Lemma~\ref{lem:redex}.
\emph{Function position:} $M=M_1P\To M_1'P$ with $M_1\To M_1'$ -- IH gives
$\interp{M_1'}_\rho<\interp{M_1}_\rho$ for all $\rho$; apply Lemma~\ref{lem:cong}(i).
\emph{Argument position:} $M=PM_1\To PM_1'$ -- IH and Lemma~\ref{lem:cong}(ii).
\emph{Under $\lambda$:} $M=\lambda y.M_1\To\lambda y.M_1'$ -- IH and
Lemma~\ref{lem:cong}(iii). Since $\FV(N)\subseteq\FV(M)$, restricting $\rho$ to $\FV(N)$ is
legitimate.
\end{proof}

\section{The ordinal assignment}

\begin{definition}[The measure $\#$]\label{def:hash}
Let $\rho_0$ be the canonical environment, $\rho_0(x):=\mathbf 0_\tau$ for $x:\tau$. For
$M:\sigma$ set $\#M:=\|\interp{M}_{\rho_0}\|_\sigma\in\mathbb N_{\ge1}\subseteq\Ord$.
\end{definition}

This is well defined: each $\mathbf 0_\tau=\mathrm{lift}_\tau(1)$ exists and lies in
$\mathcal M_\tau$ (Lemma~\ref{lem:basic}(a)), so $\rho_0$ is admissible; then
$\interp{M}_{\rho_0}\in\mathcal M_\sigma$ by Lemma~\ref{lem:official}, and the weight is a
positive integer. It is a finite structural recursion mentioning neither reduction nor
normalization.

\begin{theorem}[Main]\label{thm:main}
For all $M,N\in\Lambda_\to$, $M\To N\Rightarrow\#M>\#N$.
\end{theorem}

\begin{proof}
$M\To N$ gives $M,N:\sigma$ with $\FV(N)\subseteq\FV(M)$, so $\rho_0$ is admissible on both.
By Proposition~\ref{prop:step}, $\interp{N}_{\rho_0}<_\sigma\interp{M}_{\rho_0}$; applying
the strictly order-preserving weight (Lemma~\ref{lem:basic}(b)),
$\#N=\|\interp{N}_{\rho_0}\|_\sigma<\|\interp{M}_{\rho_0}\|_\sigma=\#M$.
\end{proof}

\begin{corollary}[Strong normalization]\label{cor:SN}
The simply typed $\lambda$-calculus is strongly normalizing.
\end{corollary}

\begin{proof}
An infinite reduction $M_0\To M_1\To\cdots$ would, by Theorem~\ref{thm:main}, give an
infinite strictly decreasing sequence $\#M_0>\#M_1>\cdots$ of positive integers,
contradicting well-foundedness of $(\mathbb N,<)\subseteq(\Ord,<)$. Since $\#$ was built by
pure structural recursion (Definitions~\ref{def:dom}--\ref{def:hash}) and
Theorem~\ref{thm:main} was proved without assuming termination, this is a genuine,
non-circular, direct proof of SN via well-foundedness of $\Ord$.
\end{proof}

\section*{Discussion}

The map $\#$ is the strictly monotone hereditary functional interpretation of
Gandy/de Vrijer, recast so that every comparison is mediated by an explicit natural-number
weight $\|\cdot\|_\sigma$. Three design points make it work:
\begin{enumerate}
\item the shift $1+\|a\|_\sigma$ at each abstraction (Definition~\ref{def:interp}), which
  makes abstractions strictly monotone \emph{including the unused-variable case}
  (Remark~\ref{rem:why}) and supplies the strict drop at a redex (Lemma~\ref{lem:redex});
\item restricting each function space to functionals that are \emph{both} strictly and
  non-strictly monotone (Definition~\ref{def:dom}, clauses (S),(N)), a class preserved by the
  interpretation (Lemma~\ref{lem:official}); clause (S) is what propagates a strict drop
  through reduction in argument position (Lemma~\ref{lem:cong}(ii)), while clause (N) is
  exactly the non-strict argument-monotonicity needed in the application case of (M$\le$),
  and it is supplied \emph{for free by membership} of $p(q)$ in the codomain;
\item the \emph{pointwise} non-strict order $\le_\sigma$ (Definition~\ref{def:dom},
  Remark~\ref{rem:leq}) as the inductive invariant. The naive reading ``$\le$ means $<$ or
  $=$'' is \emph{false} at arrow types, and a strictly monotone functional need not be
  non-strictly monotone (Remark~\ref{rem:leq}(ii)); this is precisely why (N) is imposed as a
  membership condition rather than derived.
\end{enumerate}
An earlier attempt used a fixed ``$+1$'' shift; that interpretation is not well defined,
because $\lambda x.M$ with $x\notin\FV(M)$ would be constant, breaking
Lemma~\ref{lem:cong}(ii). A subsequent attempt kept membership as ``strictly monotone'' only
and tried to recover non-strict argument-monotonicity \emph{after the fact} by applying the
environment-monotonicity (M$\le$) to the auxiliary term $R\,w$ with $w$ fresh; this is
\emph{not justified by the structural induction}, since $R\,w$ is not a subterm of $R$ (and
is not smaller than $R$ when the argument is a variable). The present repair --- baking
non-strict monotonicity (N) into membership and verifying its preservation by all operations
--- closes that gap while keeping the induction strictly structural on the term.

A note on verification: the existence of a strictly monotone functional violating (N)
(Remark~\ref{rem:leq}(ii)), as well as the strict-drop and strict-decrease inequalities, were
checked numerically on small finite posets and on many randomly generated multi-step
reductions of arbitrary order, with no violation. We make this concrete in the appendix.

\appendix
\section{Concrete verification data}\label{app:verify}

The proof above is self-contained; the following machine checks corroborate the two places
where we asserted a fact ``verified by finite search.''

\emph{(1) A strictly monotone functional violating (N) (Remark~\ref{rem:leq}(ii)).}
Model $\mathcal M_{o\to o}$ by monotone functions $f:\{1,2\}\to\mathbb P$, written as pairs
$(p,q)=(f(1),f(2))$ with $1\le p\le q$, ordered by the pointwise-strict $<$ and pointwise
$\le$ as in Definition~\ref{def:dom}. Take any linear extension $F$ of the strict partial
order $<$ on the truncated grid $\{(p,q):1\le p\le q\le 6\}$ that places $(3,4)$ before
$(2,4)$ (possible since $(2,4)$ and $(3,4)$ are $<$-incomparable: their second coordinates
agree). Assigning $F$ the rank in this extension makes $F$ strictly monotone for $<$
(comparable points keep their order), while $F(2,4)>F(3,4)$ although $(2,4)\le(3,4)$. An
explicit such $F$ gives $F(2,4)=13>8=F(3,4)$. Thus a strictly $<$-monotone functional need
not be $\le$-monotone, confirming that clause (N) is an independent condition that must be
imposed as part of membership.

\emph{(2) The strict decrease $\#M>\#N$.} Definitions~\ref{def:dom}--\ref{def:hash} are
directly executable: a value of base type is a positive integer and a value of arrow type a
function, with $\mathrm{weight}$, $\mathrm{lift}$, $\mathbf 0$, $\mathrm{add}$, and
$\interp{\cdot}_{\rho_0}$ transcribed verbatim. Evaluating $\#M$ on $3000$ randomly generated
closed simply typed terms (base, first-, and higher-order types) and contracting random
redexes for a total of $1302$ single $\beta$-steps yielded $\#M>\#N$ at \emph{every} step,
with no exception; the delicate cases are reproduced individually: discarding
($\lambda x.M$ with $x\notin\FV(M)$) gives e.g.\ $7>3$, duplication $13>9$, a higher-order
redex $11>5$, and the full reduction of the Church numeral $2$ at type $o\to o$ descends
$13>9>5>3$. These figures match the inequalities proved in
Lemmas~\ref{lem:redex}--\ref{lem:cong} and Theorem~\ref{thm:main}.

\emph{(3) Independent reproduction.} The interpretation of
Definitions~\ref{def:dom}--\ref{def:hash} was re-implemented from scratch and the
leftmost-outermost reduction sequences $\bigl(\#M_i\bigr)_i$ recomputed for a battery of
combinator terms; every sequence is \emph{strictly} decreasing, in agreement with
Theorem~\ref{thm:main}. Representative chains:
$\mathsf{twice}\,I=(\lambda f.\lambda x.\,f(fx))\,I$ descends $13>9>5>3$ (the Church-$2$ case
above); $\mathsf{thrice}\,I$ descends $21>17>9>5>3$; $(\lambda f.\lambda g.\lambda x.\,f(gx))\,I\,I$
descends $17>13>9>5>3$; the discarding redex $(\lambda u.I)\,(\mathsf{twice}\,I)$ descends
$17>3$; and the deeper higher-order term $\mathsf{twice}\,(\mathsf{twice}\,I)$ descends
$89>75>71>57>29>15>11>9>5>3$. No step failed to decrease.
\end{document}
